\chapter{Schwierigkeiten}

%Sollte vorerst passen, wird aber noch ausgebaut, mit möglichen Lösungen/unseren Lösungen, weitere Schwieriegkeiten was die Daten angeht
Bei der späteren Erfassung könnte es später zu Schwierigkeiten kommen, diese sind mitunter die verschiedene Konditionen des menschlichen Körpers sowie der Umgebung als auch das anbringen der Hardware an dem menschlichen Körper.

\section{Verschiedene Konditionen}
Obwohl das Gehen eine typisch menschliche Art der Fortbewegung ist, ist jeder Bewegungsablauf individuell sehr unterschiedlich. Diese werden von persönlichen Faktoren wie Konstitution, Beweglichkeit, Kraft, Koordinationsfähigkeit der beteiligten Muskel, die Stimmung oder sogar Umweltfaktoren wie Bodenbeschaffenheit und Umgebung sowie das Bewegungsziel beeinflusst \cite{suppe2013fbl}. Falls mehrere Personen getestet werden, muss sichergestellt werden, dass die Variabilität der Versuchenden nicht zu einer fehlerhaften Klassifikation führt. Zudem stellt sich die Frage, wie viele Personen notwendig sind, um aussagekräftige Daten zu erhalten und wie man eine Standardisierung der Testbedingungen gewährleistet.

\section{Hardware}
Eine der Hauptschwierigkeiten bei der Erfassung von Bewegungsdaten mit dem Arduino besteht in der Platzierung und dem Gewicht der Energieversorgung. Da der Arduino für mobile Anwendungen eine externe Stromquelle wie eine Powerbank benötigt, kann das zusätzliche Gewicht die natürlichen Bewegungsabläufe beeinträchtigen. Besonders beim Rennen kann die Powerbank störend wirken und die Bewegungsfreiheit einschränken. Ein erhöhtes Gewicht kann die Bewegungsfähigkeit der versuchenden erheblich beeinträchtigen \cite{monfrini2023technological} und somit zu verfälschten Messdaten führen, da die Bewegungen nicht mehr natürlich und ungehindert ausgeführt werden. Es ist daher wichtig, eine leichte und ergonomische Lösung für die Energieversorgung zu finden, um die Genauigkeit der Bewegungserfassung zu gewährleisten.

\section{Messfrequenz und Datenerfassung}
Eine zentrale Herausforderung ist die Festlegung der optimalen Messfrequenz. Eine zu niedrige Frequenz kann dazu führen, dass wichtige Bewegungsmuster nicht erfasst werden, während eine zu hohe Frequenz eine große Menge an Daten erzeugt, die schwer zu verarbeiten sind und die Systemleistung beeinträchtigen kann. Zusätzlich kann es durch Verzögerungen oder Signalstörungen zu Datenlücken kommen, die die Analyse verfälschen.